\acresetall

\chapter{Conclusion} \label{ch:60-conclusion}

\begin{foldable}
  This thesis addresses knowledge distillation along two complementary axes.
  The first is the \textit{model axis}: Chapters~\ref{ch:30-research01} and~\ref{ch:40-research02} compress a learned function under progressively tighter access restrictions---from few-shot black-box to fully data-free black-box---where the teacher withholds both its training data and its internal representations.
  The second is the \textit{data axis}: Chapter~\ref{ch:50-research03} compresses the training corpus itself, distilling a large text dataset into a compact human-readable surrogate under the additional constraint that tokens are discrete and non-differentiable.
  Together, these two axes trace a complete arc from model compression to data compression, both navigated under practical restrictions that standard \ac{KD} methods assume away.
\end{foldable}

\begin{foldable}
  Across both axes, \textit{distributional fidelity}---the degree to which the compact training signal faithfully represents the source distribution---is the binding constraint.
  Each contribution treats fidelity not as a property to hope for at sufficient scale, but as an objective to enforce by design: through adaptive confidence filtering, through structurally diverse synthesis priors, and through trajectory-aware influence weighting coupled with globally optimal transport.
  Section~\ref{sec:61-conclusion-summary} synthesises how each contribution realises this principle; Section~\ref{sec:62-conclusion-future} maps what the programme leaves open.
\end{foldable}

\section{Summary of Contributions} \label{sec:61-conclusion-summary}

\begin{foldable}
  The three contributions trace a progression of increasingly tight constraints, each answered by a more principled form of active fidelity enforcement.
  What unifies them is not a shared architecture but a shared design stance: when data is scarce, access is restricted, or tokens are discrete, fidelity to the source distribution must be built in, not assumed.
\end{foldable}

\begin{foldable}
  \textbf{DivBFKD} (Chapter~\ref{ch:30-research01}) addresses the few-shot black-box \ac{KD} problem, where the student has access to only a small unlabeled image set and may query the teacher solely for top-1 predictions.
  The key observation is that diversity of the training distribution---not merely its size---determines whether WGAN-synthesised samples can provide useful gradient signal.
  DivBFKD enforces fidelity by introducing high-confidence images, those the teacher classifies with high certainty, as real-image exemplars in the WGAN discriminator loop on-the-fly.
  Adaptive per-class confidence thresholds prevent the teacher's class-wise accuracy bias from skewing which images are admitted, ensuring that the selected exemplars represent the teacher's decision boundaries rather than its easiest cases.
  An ablation study confirms that removing adaptive thresholds---reverting to a global cutoff---degrades performance, validating that class-wise calibration of the filtering criterion is load-bearing, not incidental.
  DivBFKD achieves state-of-the-art performance among few-shot black-box \ac{KD} methods across multiple image classification datasets and architectures.
\end{foldable}

\begin{foldable}
  \textbf{DIPKD} (Chapter~\ref{ch:40-research02}) addresses the extreme case of access restriction: no real data at all, and a teacher that returns only top-1 labels.
  The fidelity problem becomes entirely generative---the student must be trained on a synthetic distribution that covers the teacher's decision surface without any real-image anchor.
  DIPKD addresses this through a three-phase pipeline where each phase serves a distinct purpose.
  The Synthesis phase constructs hierarchical multi-scale image priors from uniform noise with geometric transforms and semantic cutmixing, producing structural diversity before any teacher query.
  The Contrast phase passes these priors through a primer student backbone and applies a contrastive objective that optimises the synthetic dataset directly, pushing candidates apart in embedding space to prevent mode collapse.
  The Distillation phase then recovers soft probabilistic labels from the primer student---approximating the dark knowledge the black-box teacher withholds---and uses them alongside hard teacher labels to train the final student.
  DIPKD achieves state-of-the-art performance across twelve benchmarks, with the largest gains in fine-grained and domain-specific settings where synthetic diversity is most constrained.
\end{foldable}

\begin{foldable}
  \textbf{TAKE} (Chapter~\ref{ch:50-research03}) shifts the problem domain from image model compression to text dataset distillation~(\ac{DD}), and with it the distributional fidelity problem shifts from generation to selection.
  The goal is to replace a large labeled text corpus with a compact surrogate that preserves downstream task performance.
  TAKE is the first text \ac{DD} method to simultaneously satisfy three desiderata: knowledge-based weighting of samples by their task-aligned contribution, a dataset-level distillation objective that aligns the selected corpus with the full source distribution, and human-readable output that is auditable and transferable across architectures.
  Trajectory-aware influence scoring computes importance weights from reciprocal gradient norms accumulated across the training trajectory, correcting the hard-sample bias that convergence-point scoring introduces by assigning inflated importance to noisy or boundary examples.
  An \ac{LLM} generation pool produces diverse textual candidates, and discrete \ac{OT} via the Sinkhorn algorithm selects the subset that globally minimises the entropically regularised transport cost between the candidate pool and the source distribution.
  The importance-weighted \ac{OT} objective combines task alignment and distributional coverage as complementary objectives: importance alone directs the budget but ignores global structure; coverage alone ensures spread but ignores contribution; together they close failure modes that neither addresses individually.
  TAKE matches or exceeds prior text \ac{DD} baselines across six language benchmarks at extreme compression, establishing the first end-to-end text \ac{DD} method with theoretical guarantees.
\end{foldable}

\begin{foldable}
  Across all three contributions, the fidelity-first design stance enables deployment in regimes that scale-dependent assumptions foreclose: edge and embedded devices for DivBFKD and DIPKD, and low-resource, privacy-constrained, and carbon-sensitive language model development for TAKE.
  Model compression reduces inference cost and latency; dataset distillation reduces training expenditure, data-governance risk, and cumulative carbon footprint across full development lifecycles.
  The practical impact scales with the tightness of the constraint: the more restricted the access regime, the larger the relative gain from enforcing fidelity actively rather than assuming it.
\end{foldable}

\section{Limitations and Future Directions} \label{sec:62-conclusion-future}

\begin{foldable}
  Each contribution operates within boundaries that are not incidental but structural: they follow from the design choices that make each method tractable under its respective access constraints.
  Mapping these boundaries is therefore a prerequisite for extending the programme, not a concession to its incompleteness.
\end{foldable}

\subsection{Technical Limitations} \label{subsec:62-conclusion-future-technical}

\begin{foldable}
  DivBFKD inherits the instability of adversarial training at scale: the WGAN loop is sensitive to generator-discriminator balance, and on-the-fly injection of high-confidence real images interacts with discriminator convergence in ways we have not fully characterised at larger resolutions.
  The adaptive thresholds are calibrated against the teacher's class-wise accuracy on the available few-shot set and therefore correct for distributional bias within those classes; they are blind to domain shift between the few-shot set and the teacher's original training distribution.
  More fundamentally, high-confidence filtering guarantees teacher agreement but not task-relevant diversity: a set of confident images may still be clustered in feature space if the teacher is confidently right for consistent visual reasons.
\end{foldable}

\begin{foldable}
  DIPKD's Synthesis phase generates priors from hierarchical multi-scale uniform noise with geometric transforms and semantic cutmixing.
  This procedure yields structural diversity but is semantically agnostic: the priors may not cover the natural image manifold in fine-grained or domain-specific tasks where intra-class variation is subtle.
  The primer student is pre-selected by architecture without a criterion matching its geometric capacity to the synthetic distribution it will process, and the soft labels it produces are an approximation of the teacher's dark knowledge, with no principled bound on the gap between primer and teacher distributions.
\end{foldable}

\begin{foldable}
  TAKE's trajectory-aware influence scores are computed from reciprocal gradient norms; trajectory integration substantially mitigates hard-sample bias, but the scores lack Fisher preconditioning and rankings among samples with similar gradient magnitudes may be imperfectly calibrated.
  The discrete \ac{OT} cost matrix is defined in the encoder's embedding space, so encoder misalignment propagates directly into the transport plan.
  The \ac{LLM} generation pool fine-tuned on the source corpus risks memorizing rare examples, raising privacy concerns under membership inference, and the Sinkhorn solver scales quadratically in pool size, constraining candidate diversity at extreme compression.
  Across all three contributions, the shared cross-cutting limitation is the assumption of teacher correctness: errors, biases, and long-tail failures in the teacher propagate through each framework without detection or correction.
\end{foldable}

\subsection{Theoretical Open Problems} \label{subsec:62-conclusion-future-theoretical}

\begin{foldable}
  The fidelity-efficiency Pareto frontier is uncharacterised in all three regimes.
  No formal bound relates the distributional coverage of a synthetic or selected set to the generalization gap of a student trained on it, because the relevant bound must account jointly for the synthesis or selection procedure, the student architecture, and the task loss.
  Privacy-fidelity coupling is similarly unexplored: methods that enforce fidelity by anchoring synthesis or selection to real data---as both DivBFKD and TAKE do---may inadvertently increase leakage under membership inference, yet the privacy-fidelity trade-off has no formal treatment in the current framework.
\end{foldable}

\begin{foldable}
  The hard-sample bias that TAKE's trajectory integration corrects is empirically documented but causally unresolved: it is not known whether the bias is a selection effect, in which high-gradient samples are disproportionately noisy, or a representation effect, in which convergence-point gradients systematically underweight samples that shaped early feature learning.
  Resolving this distinction has direct implications for when and how trajectory integration should be designed.
  Finally, evaluation protocols for black-box \ac{KD} and text \ac{DD} are unstandardised across the field---different works vary teacher architectures, student capacity ratios, compression factors, and test procedures---making reliable cross-method comparison and characterisation of fidelity gains difficult.
\end{foldable}

\subsection{Future Directions} \label{subsec:62-conclusion-future-directions}

\begin{foldable}
  The most direct extension of DivBFKD is to replace the WGAN generator with a diffusion model conditioned on teacher predictions, which would simultaneously address training instability and semantic agnosticism of the synthesis prior.
  Per-class confidence gates could be learned via bilevel optimisation, with the outer loop minimising student-teacher disagreement on a held-out set, replacing the current heuristic threshold with a criterion directly tied to the distillation objective.
  Analysing the robustness of DivBFKD-distilled models to covariate shift and adversarial perturbation---as flagged in Chapter~\ref{ch:30-research01}---is a prerequisite for deployment assurance, since capability transfer and safety transfer are not guaranteed to coincide.
\end{foldable}

\begin{foldable}
  DIPKD's synthesis prior could be grounded in teacher-gradient saliency maps, even under black-box access, using gradient-free attribution methods to condition noise initialisation on the teacher's discriminative regions.
  Neural architecture search applied to primer student selection, guided by a diversity criterion on the synthetic distribution, would replace unprincipled pre-selection.
  Replacing logit matching in the Distillation phase with metric learning on primer embeddings would reduce sensitivity to the primer-teacher distribution gap by targeting structural similarity rather than output matching.
\end{foldable}

\begin{foldable}
  For TAKE, Fisher-preconditioning the influence scores is the natural next step: it would account for loss-landscape curvature and provide better-calibrated rankings, at the cost of Hessian estimation.
  Differential privacy for the \ac{LLM} generation pool under explicit $\varepsilon$ budgets would formalize the privacy guarantee; greedy lexicographic selection or neural hashing for candidate embeddings would reduce the quadratic scaling of discrete \ac{OT} and enable larger, more diverse pools.
  The broader agenda is integrating distributional fidelity with robustness, fairness, and calibration as co-equal design constraints from the outset---not as post-hoc evaluation criteria applied after the fidelity objective has already been optimised.
\end{foldable}

